\chapter{Modalit\`a di esecuzione}
\label{cap:07-modalita-esecuzione}

Una funzione registrata su \nf{} dichiara \emph{dove} vive attraverso il campo
\code{executionMode}, che ammette tre valori. La distinzione non riguarda
il protocollo di invocazione --- che \`e HTTP in tutti i casi in cui esiste un
processo separato --- ma la responsabilit\`a sul ciclo di vita del processo che
esegue il codice.

\begin{figure}[htbp]
  \centering
  \resizebox{0.95\textwidth}{!}{% nanofaas execution modes (LOCAL / POOL / DEPLOYMENT) and the cold-start vs
% warm-start latency contrast that motivates warm pools.
\documentclass[border=8pt]{standalone}
\usepackage{nanofaas-figs}
\begin{document}
\begin{tikzpicture}[nf]

  % ============ (a) dispatch targets per mode ============
  \node[core, minimum width=30mm] (disp) {PoolDispatcher};

  \node[neutral, below=13mm of disp, xshift=-52mm, minimum width=30mm] (local)
    {\textbf{LOCAL}\\{\scriptsize in-process}};
  \node[neutral, below=13mm of disp, minimum width=30mm] (pool)
    {\textbf{POOL}\\{\scriptsize warm pod pool}};
  \node[neutral, below=13mm of disp, xshift=52mm, minimum width=34mm] (dep)
    {\textbf{DEPLOYMENT}\\{\scriptsize backend provider}};

  \draw[flowc] (disp) -- (local) node[tag,midway,left=0.5mm]{test};
  \draw[flowc] (disp) -- (pool);
  \draw[flowc] (disp) -- (dep);

  % targets
  \node[backend, below=9mm of local, minimum width=30mm] (localt)
    {handler in JVM\\{\scriptsize no container}};
  \node[backend, below=9mm of pool, minimum width=34mm] (poolt)
    {pre-warmed pods\\{\scriptsize OpenWhisk-style}};
  \node[backend, below=9mm of dep, xshift=-11mm, minimum width=26mm] (k8s)
    {Kubernetes\\{\scriptsize Deploy+Svc}};
  \node[backend, below=9mm of dep, xshift=19mm, minimum width=26mm] (clocal)
    {container-local\\{\scriptsize Docker CLI}};

  \draw[flowc] (local) -- (localt);
  \draw[flowc] (pool) -- (poolt);
  \draw[flowc] (dep.south) -- (k8s.north);
  \draw[flowc] (dep.south) -- (clocal.north);

  % ============ (b) cold-start vs warm-start timeline ============
  \newcommand{\phase}[6]{% x1,x2,y,fill,draw,label
    \draw[fill=#4, draw=#5, line width=0.5pt, rounded corners=1pt]
      (#1,#3-0.3) rectangle (#2,#3+0.3);
    \node[font=\sffamily\scriptsize, text=nfink] at ({(#1+#2)/2},#3) {#6};}

  \coordinate (t0) at (-5.2,-7.9);   % origin of the timeline block
  \begin{scope}[shift={(t0)}]
    \node[clabel, anchor=west] at (-0.2,1.95)
      {Cold start vs.\ warm start \textnormal{\footnotesize— first request latency}};

    % axis
    \draw[flowc] (0,-1.5) -- (11.2,-1.5) node[tag, right]{time};

    % cold track (DEPLOYMENT, first hit)
    \node[tag, anchor=east, align=right] at (-0.15,0.35)
      {DEPLOYMENT\\(cold)};
    \phase{0}{1.4}{0.35}{nfslatebg}{nfslate}{schedule}
    \phase{1.4}{5.0}{0.35}{nfamberbg}{nfamber}{pull image}
    \phase{5.0}{6.8}{0.35}{nfgreenbg}{nfgreen}{start ctr}
    \phase{6.8}{9.0}{0.35}{nfpurplebg}{nfpurple}{app ready}
    \phase{9.0}{10.6}{0.35}{nfbluebg}{nfblue}{serve}

    % warm track (POOL / warm pool)
    \node[tag, anchor=east, align=right] at (-0.15,-0.65)
      {POOL /\\warm pool};
    \phase{0}{0.7}{-0.65}{nfgreenbg}{nfgreen}{ready}
    \phase{0.7}{2.3}{-0.65}{nfbluebg}{nfblue}{serve}

    % cold-start penalty bracket
    \draw[decorate, decoration={brace, amplitude=4pt, raise=1pt}, draw=nfred]
      (9.0,0.72) -- (0.0,0.72);
    \node[tag, text=nfred] at (4.5,1.15) {cold-start penalty};
    \draw[dashed, draw=nfslate] (0.7,-1.0) -- (0.7,-1.5);
    \draw[dashed, draw=nfslate] (9.0,0.0) -- (9.0,-1.5);
  \end{scope}

\end{tikzpicture}
\end{document}
}
  \caption{Le tre modalit\`a di esecuzione e il contrasto fra avvio a freddo e a
  caldo.}
  \label{fig:modalita}
\end{figure}

\section{\code{LOCAL}: esecuzione in-process}

\code{LocalDispatcher} \`e sei righe di codice:

\begin{lstlisting}[language=Java]
@Component
public class LocalDispatcher implements Dispatcher {
    @Override
    public CompletableFuture<DispatchResult> dispatch(InvocationTask task) {
        return CompletableFuture.completedFuture(
                DispatchResult.warm(InvocationResult.success(task.request().input())));
    }
}
\end{lstlisting}

Restituisce l'input come output, immediatamente, dichiarandosi caldo. \`E un eco.

Il suo scopo \`e sperimentale, non applicativo. Fornisce un percorso di
invocazione la cui componente di esecuzione ha costo nullo e varianza nulla, e
quindi rende misurabile il costo del control plane in isolamento: qualunque
latenza osservata su una funzione \code{LOCAL} \`e attribuibile interamente alla
piattaforma. \`E anche il modo in cui la maggior parte dei test di integrazione
esercita il percorso completo senza avviare container.

\section{\code{EXTERNAL}: inoltro a un endpoint esistente}

In modalit\`a \code{EXTERNAL} la funzione \`e ospitata fuori dal control plane,
che ne conosce soltanto l'URL (\code{endpointUrl}) e vi inoltra le invocazioni.
Non ne gestisce n\'e la creazione, n\'e le repliche, n\'e la distruzione. \`E,
letteralmente, un passthrough.

Il codice ne trae due conseguenze. La prima \`e che l'assenza di
\code{endpointUrl} \`e un errore rilevato al momento del dispatch, non della
registrazione:

\begin{lstlisting}[language=Java]
String endpoint = task.functionSpec().endpointUrl();
if (endpoint == null || endpoint.isBlank()) {
    return CompletableFuture.completedFuture(DispatchResult.warm(
            InvocationResult.error("EXTERNAL_ENDPOINT_MISSING",
                                   "endpointUrl is required for EXTERNAL mode")));
}
\end{lstlisting}

La seconda \`e che il blocco di scaling non viene risolto per le funzioni
\code{EXTERNAL} (\cref{cap:06-nucleo}): il control plane non pu\`o
ragionevolmente decidere il numero di repliche di un deployment che non
possiede.

La modalit\`a \code{EXTERNAL} \`e quella che il servizio di riferimento
\code{warm-echo} usa, ed \`e il veicolo principale degli esperimenti in cui il
runtime della funzione dev'essere tenuto costante mentre si variano le politiche
del control plane.

\section{\code{DEPLOYMENT}: intento gestito}

\code{DEPLOYMENT} \`e la modalit\`a completa: il control plane si assume la
responsabilit\`a di far esistere le repliche. \`E anche il default in assenza di
dichiarazione esplicita.

La sua realizzazione segue un principio: \code{DEPLOYMENT} \`e un
\emph{intento}, non un meccanismo. Il control plane non sa creare deployment;
sa dichiarare che ne serve uno e delegarne la realizzazione a un provider.

\subsection{L'interfaccia del provider}

\begin{lstlisting}[language=Java]
public interface ManagedDeploymentProvider {
    String backendId();
    boolean isAvailable();
    boolean supports(FunctionSpec spec);
    ProvisionResult provision(FunctionSpec spec);
    void deprovision(String functionName);
    void setReplicas(String functionName, int replicas);
    int getReadyReplicas(String functionName);
    default ReplicaStatus getReplicaStatus(String functionName) { ... }
}
\end{lstlisting}

L'interfaccia \`e piccola e i suoi primi tre metodi sono i pi\`u interessanti,
perch\'e non riguardano il provisioning ma la \emph{selezione}. Un provider deve
poter dichiarare di esistere (\code{backendId}), di essere utilizzabile in questo
momento (\code{isAvailable} --- il provider Kubernetes non lo \`e senza un
kubeconfig valido) e di poter servire una specifica particolare
(\code{supports}).

Il metodo \code{provision} restituisce un \code{ProvisionResult} che contiene
l'endpoint risultante e i metadati degli oggetti creati. Il ritorno di un
endpoint anzich\'e di un handle opaco \`e ci\`o che consente al dispatcher di
trattare uniformemente \code{DEPLOYMENT} ed \code{EXTERNAL} dopo il
provisioning: una volta che l'endpoint esiste, invocarlo \`e la stessa
operazione.

\subsection{La risoluzione del backend}

\code{DeploymentProviderResolver} sceglie il provider con una precedenza a tre
livelli: suggerimento esplicito per registrazione, default di piattaforma
(\code{nanofaas.deployment.default-backend}), e infine risoluzione implicita.

La risoluzione implicita \`e la parte in cui il progetto compie una scelta
difendibile ma non ovvia:

\begin{lstlisting}[language=Java]
if (candidates.isEmpty()) {
    throw new IllegalStateException("No managed deployment provider available for function ...");
}
if (candidates.size() > 1) {
    throw new IllegalStateException("Ambiguous managed deployment provider selection for function '"
            + spec.name() + "': " + ids);
}
return candidates.getFirst();
\end{lstlisting}

Con due provider entrambi disponibili e capaci, il resolver \textbf{fallisce}
anzich\'e scegliere. Una piattaforma con un ordine di preferenza implicito ---
``prova Kubernetes, altrimenti Docker'' --- funzionerebbe, e produrrebbe
deployment su un backend diverso da quello atteso ogni volta che l'ambiente
cambia sotto i piedi dell'operatore. Un'ambiguit\`a dichiarata \`e preferibile a
una preferenza nascosta. L'operatore risolve fissando
\code{default-backend}, che \`e un atto esplicito.

\subsection{Degradazione da \code{DEPLOYMENT} a \code{EXTERNAL}}

Esiste un caso in cui il provisioning fallisce ma la registrazione riesce:

\begin{lstlisting}[language=Java]
public ProvisionResult resolveAndProvision(FunctionSpec spec, String backendHint) {
    try {
        return resolve(spec, backendHint).provision(spec);
    } catch (IllegalStateException ex) {
        if (spec.endpointUrl() != null && !spec.endpointUrl().isBlank()) {
            return new ProvisionResult(spec.endpointUrl(), null,
                    ExecutionMode.EXTERNAL, ex.getMessage(), Map.of());
        }
        throw ex;
    }
}
\end{lstlisting}

Se la funzione dichiara \emph{anche} un \code{endpointUrl}, un fallimento del
provisioning la degrada a \code{EXTERNAL} anzich\'e respingerla. \`E la
degradazione che consente allo stesso manifesto di funzione di essere usato su
un cluster Kubernetes e su una macchina di sviluppo priva di backend gestito.

La degradazione \`e \emph{osservabile}, ed \`e questo a renderla accettabile
secondo il criterio del \cref{cap:03-requisiti-principi}. Il tipo
\code{DeploymentMetadata} conserva contemporaneamente la modalit\`a richiesta,
quella effettiva e la ragione della differenza:

\begin{lstlisting}[language=Java]
public record DeploymentMetadata(
        ExecutionMode requestedExecutionMode,
        ExecutionMode effectiveExecutionMode,
        String deploymentBackend,
        String degradationReason,
        String effectiveEndpointUrl,
        Map<String, String> deploymentObjects
) { ... }
\end{lstlisting}

L'API di ispezione della funzione restituisce entrambe le modalit\`a e la
ragione, cos\`i che un operatore che si aspettava un deployment Kubernetes possa
scoprire di averne ottenuto un passthrough --- e perch\'e --- senza leggere i
log.

\section{Risveglio da zero repliche}

Con un deployment scalato a zero, un'invocazione arriva prima che esista una
replica capace di servirla. Il nucleo affronta il caso con
\code{DeploymentWakeUpGate}: prima del dispatch verso l'endpoint, se la funzione
\`e gestita e non ha repliche pronte, la richiesta attende che ne compaia una.

Il meccanismo \`e coordinato: pi\`u invocazioni concorrenti sulla stessa funzione
condividono la stessa attesa anzich\'e emettere ciascuna la propria richiesta di
scale-up. \`E il minimo indispensabile per evitare che una raffica su una
funzione fredda si traduca in una raffica di comandi al backend.

\begin{damisurare}
La distribuzione del tempo di risveglio da zero repliche \`e strumentata
(\code{function\_cold\_start\_ms}) ma non \`e stata caratterizzata
sistematicamente: le campagne di misura del \cref{cap:23-prestazioni-base}
partono da repliche gi\`a calde. \todomis{distribuzione del tempo di risveglio
per backend Kubernetes e container locale}
\end{damisurare}

\section{Confronto}

\begin{table}[htbp]
\centering
\caption{Le tre modalit\`a a confronto.}
\label{tab:modalita}
\small
\begin{tabularx}{\textwidth}{@{}lXXX@{}}
\toprule
 & \textbf{\code{LOCAL}} & \textbf{\code{EXTERNAL}} & \textbf{\code{DEPLOYMENT}} \\
\midrule
Processo funzione & nessuno & preesistente & creato dal control plane \\
Ciclo di vita & --- & del proprietario & del control plane \\
Scaling repliche & --- & non applicabile & \code{autoscaler} o backend \\
Cold start & assente & fuori dal controllo & osservato e misurato \\
Validazione immagine & no & no & s\`i, dal validatore del backend \\
Uso previsto & misura del control plane, test & runtime fissato negli esperimenti & carichi realistici \\
\bottomrule
\end{tabularx}
\end{table}

La riga ``uso previsto'' riassume la logica del progetto: le tre modalit\`a non
sono tre livelli di completezza ma tre strumenti sperimentali. \code{LOCAL}
isola il control plane, \code{EXTERNAL} lo mette a confronto con un runtime che
non cambia fra un esperimento e l'altro, \code{DEPLOYMENT} chiude il ciclo e
include l'infrastruttura nella misura.
